\subsection{Performance evaluation}

\begin{qfmethod}[Relative performance evaluation: total returns, market-adjusted and risk-adjusted returns, market timing, and smoothed returns]\label{meth:relative-performance-evaluation}
\qffield{Statement} Portfolio performance depends on management and on every factor that moves the prices of its holdings, most of which nobody can foresee, so good results may reflect skill or luck \citep[chap.~``Performance Evaluation and Prediction'', sec.~22.1, pp.~443--444]{alpha-harris}. Relative evaluations are therefore more valuable than absolute ones: an absolute evaluation computes the total return; a relative one compares it with a benchmark chosen to show how the portfolio would have done without active management (sec.~22.2, p.~445). Harris's measures (secs.~22.2.1--22.2.2.3, pp.~445--450):
\begin{itemize}
\item \emph{total return} over an interval, including income and paying out commissions, fees and margin flows; with capital added or withdrawn, either the \emph{internal rate of return}, approximately a time- and value-weighted average of the returns between flows, or the \emph{holding-period return} of one unit invested at the start with distributions reinvested, which differ when flows are not pro rata;
\item \emph{market-adjusted return}, the portfolio return minus a market index return;
\item \emph{risk-adjusted excess return} or realised alpha, the portfolio return minus the portfolio's beta times the market return, with the beta measured concurrently because it changes when the manager trades;
\item \emph{market-timing return}, beta times the market return minus the market return; so that
\end{itemize}
\[ r^p=\underbrace{(r^p-\beta r^m)}_{\text{excess return}}+\underbrace{(\beta r^m-r^m)}_{\text{market timing}}+\underbrace{r^m}_{\text{market}} \]
(p.~450). Harris notes that the decomposition generalises to other risk factors, such as size and growth.

\qffield{Conditions} Timeless (methods); presentation standards named by Harris (AIMR-PPS) as of 2002. Returns of illiquid assets depend on estimated values; managers who adjust them slowly, especially downwards, report smooth returns with positive serial correlation; the low volatility inflates $t$-ratios, and forced revaluation later produces large losses, an artificial peso problem \citep[same chap., box ``Smooth Sailing, Rapids Ahead'', p.~448, and sec.~22.5.2.1, pp.~469--470]{alpha-harris}. Holding-period returns compound: $+100$ percent then $-50$ percent is zero, not $+50$ percent (box, p.~447).

\qffield{Practice} Interpretation: transaction cost analysis applies the same idea over an order's horizon, subtracting beta times the index move from the arrival cost (\hyperref[meth:transaction-cost-analysis]{transaction cost analysis}); in both, the residual is not proof of skill.

\qffield{Pitfalls} Interpretation: Harris's realised alpha $r^p-\beta r^m$ omits the risk-free rate; the regression form $r^p_t-r^f_t=\alpha+\beta(r^m_t-r^f_t)+u_t$ (Jensen's alpha) differs from it by $(1-\beta)r^f$, for example $-0.9$ percentage points at $\beta=1.3$ and $r^f=3$ percent; the regression also estimates a constant beta, whereas Harris uses concurrent betas. The decomposition $r^p=r^m+(r^p-\beta r^m)+(\beta-1)r^m$ is Harris's identity. Harris calls the IRR approximately time- and value-weighted, not money-weighted, and does not discuss who controls flows; fill quality belongs to transaction cost analysis. Return smoothing is in Harris a problem of fraudulent or careless valuation; for a formal treatment see \nolinkurl{wide-getmansky-lo-makarov}.

\qffield{See also} \hyperref[meth:transaction-cost-analysis]{transaction cost analysis} (the execution analogue).

\end{qfmethod}

\begin{qfmethod}[Testing skill against luck: the $t$-test of mean excess return, its power and data needs, sample length against sampling frequency, distributional shape, smoothed and fraudulent returns, sample selection and regression to the mean]\label{meth:skill-luck-significance-tests}
\qffield{Statement} To judge whether a manager or strategy beats the market systematically, Harris compares the average market-adjusted return with ``the standard error of the mean'', the size of average return luck alone would produce; if their ratio, Student's $t$, is large, ``luck alone probably cannot explain the returns''; for an unskilled manager the one-sided $t$ exceeds 1.64 only 5 percent of the time \citep[chap.~22, sec.~22.4.1, p.~453]{alpha-harris}. With excess returns $x_1,\dots,x_n$, $t=\bar x/(s_x/\sqrt n)$. The test's \emph{power}, the probability of calling a skilled manager skilled, rises as the confidence level falls, with skill, with more years of data, and falls as luck matters more (sec.~22.4.2, p.~454).

Harris's calibration: skill of 2 percent a year, and a market-adjusted standard deviation of 7 percent from $\sqrt{16^2+14.5^2-2\times0.9\times16\times14.5}=6.97$ (p.~455); a 95 percent test then identifies the skilled manager with probability 15 percent from five years of monthly data and 23 percent from ten (p.~455; Table~22-2, p.~456), and a test with 75 percent confidence and 75 percent power needs 22 years of data, one with 95 and 95 percent 132 years (Table~22-3, p.~457). ``In practice, more than 20 years of returns data are typically required to obtain useful results for a given investment manager'' (p.~453). The cause is ``a low signal to noise ratio'' (box, p.~463).

\qffield{Notation} $x_i$ excess (market-adjusted) return in period $i$; $\mu$ and $\sigma$ its true mean and standard deviation; $\bar x$, $s_x$ their sample estimates; $n$ the number of periods; $\alpha$ the significance level (one minus Harris's confidence level) and $1-\beta$ the power; $z_q$ the standard normal $q$-quantile.

\qffield{Conditions} \emph{Power and sample size} (derived). With independent normal returns the $t$-statistic is centred near $\mu\sqrt n/\sigma$, so a one-sided test at level $\alpha$ with power $1-\beta$ needs about $n=\bigl((z_{1-\alpha}+z_{1-\beta})\,\sigma/\mu\bigr)^2$ periods: of order $(\sigma/\mu)^2$, as the proposition says, with a factor set by the error rates. In years, this reproduces Harris's tables: power 0.157 and 0.229 for five and ten years; 22.3, 132.6 and 33.1 years for (75\%, 75\%), (95\%, 95\%) and (50\%, 95\%) at 2 percent skill, and 33.1 years for (95\%, 95\%) at 4 percent. The proof's ``expected $t$-statistic is approximately $\mu\sqrt n/\sigma$'' is a large-sample statement: with ten annual observations the mean simulated $t$ is 1.00 against 0.90.

\emph{Sample length, not frequency.} ``statistical power depends primarily on the length of the sample period and not on the frequency of sampling within that period'', because total performance over a period does not depend on how often it is observed; more frequent sampling can estimate the standard error better and helps mainly when there would otherwise be fewer than 20 observations (sec.~22.4.6.1, p.~466). Simulated over ten years, the mean $t$ is 0.91 from monthly and 0.90 from daily data. With serially correlated returns the naive standard error is too small: for AR(1) returns with coefficient 0.3 the naive $t$ of a zero-mean series has standard deviation 1.36 instead of 1 (Interpretation). More observations within the same period do not add history.

\emph{Distributional shape.} The test assumes normal returns; fat tails make it less powerful, asymmetric strategies such as selling volatility (the peso problem) look skilled until the rare loss, and smoothed valuations of illiquid assets lower measured volatility, giving ``artificially high $t$-ratios, and therefore overconfident conclusions''; pyramid schemes look skilled until they collapse (sec.~22.5, pp.~467--470). A stream earning 1 percent a month with a 2 percent monthly chance of losing 40 percent has a true mean of 0.18 percent a month and shows no loss in 60 months with probability 0.30.

\qffield{Practice} \emph{Selection.} ``The sample selection bias arises when some process selects the information that you see''; fund families advertise their best funds and merge away their losers (survivorship bias), and newsletters with poor records disappear (sec.~22.6, pp.~470--473). Significance must be judged against all the managers who might have been lucky: a $t$ of 4.9 for Buffett's 36 years is extreme for one manager, but Harris estimates that the best of 10{,}000 unskilled managers would beat the market by 11.8 percent a year with probability 0.5 percent, and the best of 100{,}000 with 5 percent (box, p.~463). With normal 36-year averages, these are 0.37 and 3.6 percent; Harris's method is not stated.

\emph{Regression to the mean}: selection can raise an average of past returns but ``cannot affect subsequent returns'', so ``When selection has caused the past return average to be high, the subsequent return average is invariably lower'' (sec.~22.6.3, p.~474). In a simulation where skill has variance $\tau^2=0.25$ and luck $\sigma^2=1$, the top decile's next-period mean is 0.20 of its selection-period mean, $\tau^2/(\tau^2+\sigma^2)$ (see \hyperref[res:zero-sum-skill-luck]{zero-sum trading, skill and luck}).

\qffield{Pitfalls} Additional sources of bias include \emph{look-ahead} (timestamps, corporate actions or benchmark constituents that use future information), which belongs to \hyperref[data:trading-reference-data]{point-in-time reference data} and \hyperref[pit:label-weight-leakage]{time-validity of labels}, and \emph{fillability} (simulated fills assume unreachable liquidity), which is \hyperref[pit:backtest-cost-short-assumptions]{backtest cost assumptions}; multiple testing is Harris's best-of-many problem, whose corrections (deflated Sharpe ratios, stepdown tests) are in \hyperref[meth:deflated-sharpe-ratio]{deflated Sharpe ratio} and \hyperref[meth:multiple-testing-control]{multiple-testing control}. That most backtest errors are sampling errors before they are modelling errors is a generalisation; Harris says sample selection biases ``may be responsible for more trading losses than any other cause'' (p.~472). The calibration (2 percent skill, 7 percent noise) is Harris's assumption for diversified U.S. equity managers as of 2002, and the test is one-sided on normal monthly returns. See \nolinkurl{wide-fama-french-luck}, \nolinkurl{wide-barras-scaillet-wermers} and \nolinkurl{wide-getmansky-lo-makarov} on luck in fund returns, smoothing and hedge-fund returns.

\qffield{See also} \hyperref[meth:relative-performance-evaluation]{relative performance evaluation} (market-adjusted returns and smoothing); \hyperref[res:zero-sum-skill-luck]{zero-sum trading, skill and luck}; \hyperref[pit:parameter-overfitting]{parameter overfitting}; \hyperref[meth:backtest-evaluation-diagnostics]{backtest evaluation diagnostics}; \hyperref[meth:venue-market-share]{venue market share} (proportion intervals).

\qffield{Further reading} Harris, chap.~22, secs.~22.4--22.6.

\end{qfmethod}

\begin{qfmethod}[A GIPS composite report: which portfolios form a composite, the statistics it must present and its first required disclosures (GIPS for Firms, 2020 edition)]\label{meth:gips-composite-presentation}
\qffield{Statement} The 2020 edition of the GIPS standards for firms, effective 1 January 2020 \citep[Introduction, PDF p.~17]{wide-gips-2020}, sets composite construction in section~3 \citep[provisions 3.A.1--3.A.11, pp.~18--19]{wide-gips-2020}:
\begin{itemize}
\item a firm ``must create composites'' for strategies managed or offered as segregated accounts, and ``all actual, fee-paying, discretionary segregated accounts must be included in at least one composite''; non-discretionary portfolios must not be;
\item composites are ``defined according to investment mandate, objective, or strategy'' and include every portfolio meeting the definition, whatever its legal structure, and a change of definition ``must not be applied retroactively'';
\item new portfolios enter ``on a timely and consistent composite-specific basis''; only portfolios managed for the full measurement period count; terminated portfolios stay in history ``up to the last full measurement period''; portfolios move between composites only on documented client-directed changes or a redefinition, and their history stays with the original composite.
\end{itemize}
A composite time-weighted return report must present (4.A.1, pp.~21--22):
\begin{itemize}
\item at least five years of compliant performance, or since inception, building up to ten;
\item annual composite returns and ``the total return for the benchmark for each annual period and for all other periods for which composite returns are presented, unless the firm determines there is no appropriate benchmark'' (4.A.1(b) and~(e), PDF p.~39, printed p.~21);
\item the number of portfolios, if more than five, composite assets and total firm assets at each annual period end;
\item a measure of internal dispersion of portfolio returns, if more than five portfolios were held for the full year;
\item where monthly returns are available, the three-year annualised ex post standard deviation of composite and benchmark, from monthly returns.
\end{itemize}
It must label the periods and whether returns are gross or net of fees (4.A.3). Its disclosures begin with one of three prescribed compliance statements: verified, verified with a performance examination, or not verified. Then come the trademark statement, the definition of the firm, the composite description and the benchmark description (4.C.1--4.C.5, pp.~25--27).

\emph{No appropriate benchmark.} The benchmark column has its own exception, 4.A.1(e) above. When it is used, ``If the firm determines no appropriate benchmark for the composite exists, the firm must disclose why no benchmark is presented'' (4.C.31, PDF p.~47, printed p.~29). Money-weighted composite reports and pooled-fund reports carry the same disclosure (5.C.30, 6.C.26, 7.C.26), and the time-weighted pooled-fund report the same exception (6.A.1(e)). In GIPS advertisements the explanation is only recommended (8.D.4, 8.F.4, 8.H.1: ``should disclose''). Money-weighted composite reports, for example for closed-end funds, present since-inception money-weighted returns instead (5.A.1, p.~33).

\qffield{Notation} Composite: an aggregation of portfolios managed to a similar mandate; internal dispersion: the spread of annual returns across the composite's portfolios (the firm discloses which measure); ex post standard deviation: of 36 monthly returns, annualised.

\qffield{Conditions} The provisions are requirements of a voluntary standard: they bind a firm that claims compliance, and the firm must disclose when the three-year standard deviation is not presented because 36 monthly returns are not available (4.C, p.~30). Footnotes date some requirements: total firm assets from periods ending 31 December 2020 (n.~31), part-period returns for composites starting from 1 January 2011 (nn.~29--30).

\qffield{Practice} Interpretation: the composite rules are survivorship and cherry-picking controls. Every qualifying account must be included, terminated accounts stay in the history, nothing is applied retroactively and portfolios move only for documented reasons. The required statistics let a reader see scale, dispersion across clients and risk next to return. A toy report needs the five-year table with these columns and the compliance statement wording, not only returns.

\emph{Presenting and updating.} A firm must give the report to every prospective client when they become one and update it at least every 12 months while they remain one (1.A.11--1.A.12, p.~2), with information through the most recent annual period end within 12 months of that period end (1.A.16). If a composite loses all its portfolios its track record ends; a restarted record is presented separately, the break shown and not linked across (4.A.5, p.~22). Theoretical performance must not be linked to actual performance (1.A.27, p.~4).

\emph{Portability of past performance.} Performance from a past firm or affiliation may represent and be linked to the record of the new or acquiring firm only if, composite by composite: ``Substantially all of the investment decision makers'' are employed by the new firm; the decision-making process remains ``substantially intact and independent''; the new firm has records to support the performance; and there is no break in the track record. Otherwise it must not be linked (1.A.32, p.~5). Where the only failure is a break in the record, the performance before and after the break may be presented separately but not linked (1.A.33). Acquired non-compliant assets must comply within one year (1.A.34). The report must disclose if past-firm performance is presented and for which periods (4.C.37, p.~30). The Handbook adds that ``composite'' means the entire composite at the past firm: a record rebuilt from a subset of portfolios may appear in a GIPS Report only as supplemental information, never linked, and performance must never be presented without supporting records \citep[Provision 1.A.32, pp.~50--52]{gips-handbook-firms-2020}. An educational report is in \hyperref[ex:gips-composite-report]{the educational composite report}.

\qffield{Pitfalls} Computing conventions for internal dispersion and the 36-month figure ($n$ or $n-1$, full-year portfolios only) are in \hyperref[meth:gips-return-calculations]{GIPS calculations}. Pooled-fund reports (section~6 onwards) differ and were only searched for the matching provisions. Linking a departing team's subset of accounts as if it were the old composite.

\qffield{See also} \hyperref[def:return-measures]{return measures} (time- and money-weighted returns); \hyperref[meth:relative-performance-evaluation]{relative performance evaluation}; \hyperref[meth:skill-luck-significance-tests]{tests of skill against luck}.

\qffield{Further reading} GIPS Guidance Statements and the Q\&A database.

\end{qfmethod}
